\chapter{The probability ideas we need}
\label{chap:probability}

\section{Probability describes a chance}

A \textbf{probability} is a number between zero and one. Zero means that an
outcome is impossible under the stated rules. One means that it is certain.
A probability of \(0.25\) is the same as one quarter or twenty-five percent.

An \textbf{event}, in probability language, is a specified outcome or set of
outcomes. For example, ``the selected individual is A'' is an event. The
word event will also refer to one reproduction and replacement attempt in
the Moran model. The surrounding sentence will tell us which meaning is used.

We write \(\Pr(A)\) for the probability of a probability event called
\(A\). Here \(A\) is a temporary name for an event, not necessarily the
genetic type A. For a concrete example, we can avoid a letter and write
\(\Pr(\text{selected individual is A})\).

Probabilities concern the rules and information available before a draw.
Once a particular draw has happened, its outcome is known. The fact that an
unlikely outcome occurred does not mean its original probability was wrong.

\section{Fractions and equal chances}

Suppose a population contains five individuals. Two are A and three are B.
If each individual is equally likely to be selected, there are five equally
likely possibilities. Two of them meet the description ``selected individual
is A''. Therefore
\[
  \Pr(\text{selected individual is A})=\frac{2}{5}=0.4.
\]

The rule ``favourable possibilities divided by all possibilities'' works here
because the individual possibilities have equal chances. We must not use it
without checking that condition. Later, choosing a parent will involve
unequal reproductive weights. Choosing an individual for replacement will
still give equal chances to all individuals.

\begin{workedexample}
Think of five equally sized tickets in a bag. Two tickets say A and three
say B. Drawing one ticket uniformly means that each ticket has probability
\(1/5\). The total probability of drawing A is
\(1/5+1/5=2/5\). We add the probabilities because the two A tickets are
different possible outcomes of the same draw.
\end{workedexample}

\section{Opposite events and the total probability}

If there are only A and B individuals, a selected individual must be one
or the other. It cannot be neither. It also cannot be both in our haploid
model. Their probabilities therefore add to one:
\[
 \Pr(\text{A selected})+\Pr(\text{B selected})=1.
\]
If the first probability is \(0.4\), the second must be \(1-0.4=0.6\).

The event ``not A'' is called the \textbf{complement} of ``A''. In general,
the probability of a complement is one minus the probability of the original
event. We will use this fact when calculating the chance that the allele
count stays unchanged.

\section{Conditional probability: state what is already known}

The phrase \textbf{conditional probability} means a probability calculated
given some specified information. A vertical bar is read as ``given''.
For example,
\[
 \Pr(\text{A selected}\mid i=2,N=5)
\]
means ``the probability that A is selected, given that there are two A
individuals in a population of five''. The symbols \(i\) and \(N\) have
the meanings introduced in Chapter~\ref{chap:foundations}.

The information matters. If there were no A individuals, the chance would
be zero. If all five individuals were A, the chance would be one.

For a more abstract statement, let \(A\) and \(B\) name two probability
events. The expression \(\Pr(B\mid A)\) is the probability that \(B\)
occurs when we know that \(A\) has occurred. If \(\Pr(A)>0\), its definition is
\[
  \Pr(B\mid A)=\frac{\Pr(A\text{ and }B)}{\Pr(A)}.
\]
The numerator is the probability that both events occur. Dividing by
\(\Pr(A)\) restricts our attention to occasions on which \(A\) occurred.

We will not need to manipulate this definition repeatedly. The important
habit is to ask: \emph{what information is this probability conditioned on?}

\section{Why we multiply probabilities}

Rearranging the definition above gives the general multiplication rule:
\[
 \Pr(A\text{ and }B)=\Pr(A)\Pr(B\mid A).
\]
The expression on the right means the chance of the first event, multiplied
by the chance of the second event given the first.

Two events are \textbf{independent} when knowing that the first occurred does
not change the probability of the second. In that case
\(\Pr(B\mid A)=\Pr(B)\), so the multiplication rule becomes
\[
 \Pr(A\text{ and }B)=\Pr(A)\Pr(B).
\]

\begin{workedexample}
Suppose the chance of choosing an A parent is \(0.5\). Suppose an independent
replacement draw has probability \(0.6\) of selecting a B individual.
The chance of choosing an A parent \emph{and} replacing a B individual is
\[
  0.5\times 0.6=0.3.
\]
As an intuitive count, imagine one hundred repetitions. About fifty choose
an A parent. Of those fifty, about sixty percent, or thirty, replace B.
Thirty out of one hundred is \(0.3\). These are expected counts, not
promises about a particular batch of one hundred repetitions.
\end{workedexample}

In our model, the parent and replacement draws are independent
\emph{given the current population state}. They both use the same current
counts. We do not remove the parent before choosing a replacement. We also
do not rule out selecting the same individual twice.

This conditional statement is precise. If we mixed observations from many
different population states without recording those states, the apparent
relationship between the draws could be different.

\section{Why we add probabilities}

Events are \textbf{mutually exclusive} when they cannot both happen in the
same trial. The outcomes ``the A count goes up'', ``the A count goes down'',
and ``the A count stays the same'' are mutually exclusive.

When events are mutually exclusive, we find the probability that one or
another occurs by adding their probabilities. For example, suppose an event
has probability \(0.3\) of increasing the count and probability \(0.2\)
of decreasing it. Then
\[
 \Pr(\text{count changes})=0.3+0.2=0.5.
\]
The remaining probability is \(1-0.5=0.5\), so the count stays unchanged
with probability \(0.5\).

Do not add probabilities in this simple way when the events overlap. If
the same possible outcome belongs to both events, adding would count it
twice. The three count outcomes above do not overlap.

\section{A random variable is a numerical outcome}

A \textbf{random variable} is a numerical quantity determined by a random
outcome. Let \(Y\) denote the change in the A count during one event. The
possible values of \(Y\) are
\[
  Y=1,\qquad Y=-1,\qquad\text{or}\qquad Y=0.
\]
The value \(1\) means one more A individual. The value \(-1\) means one
fewer A individual. The value \(0\) means no change.

The symbol names the quantity before we know which value the random event
will produce. A realised value is the actual value produced in one trial.

\section{Expected value: a probability-weighted average}

The \textbf{expected value} of a random variable is its average when each
possible value is weighted by its probability. We write \(\mathbb{E}[Y]\)
for the expected value of \(Y\). The stylised letter \(\mathbb{E}\) is
read as ``expectation''.

Suppose \(Y\) equals \(1\) with probability \(0.3\), equals \(-1\) with
probability \(0.2\), and equals \(0\) with probability \(0.5\). Multiply
each value by its probability and add:
\begin{align*}
 \mathbb{E}[Y]
 &=1(0.3)+(-1)(0.2)+0(0.5)\\
 &=0.3-0.2+0\\
 &=0.1.
\end{align*}

The expected change is one tenth of an individual per event. An actual
event cannot add one tenth of an individual. It adds one, removes one, or
changes nothing. An average does not have to be a possible individual outcome.

For intuition, imagine ten events that each begin from the \emph{same}
population state. On average across many such batches, three events increase
the count and two decrease it. The average total change is one, so the
average change per event is \(1/10=0.1\).

This thought experiment resets the starting state before each event. In a
single evolving population, a changed count changes the probabilities of
later events. We cannot assume that ten successive events in one trajectory
all use the original probabilities.

\subsection{Expected change in frequency}

The frequency is the count divided by the fixed population size \(N\).
Let \(\Delta x\) mean the change in frequency during one event. The Greek
capital letter \(\Delta\) is read as ``change in''. Since the count change
is \(Y\), the frequency change is \(\Delta x=Y/N\).

For the example above with \(N=5\), the three possible frequency changes
are \(1/5\), \(-1/5\), and zero. Their expected value is
\[
 \mathbb{E}[\Delta x]=\frac{0.1}{5}=0.02.
\]
Again, a positive average does not rule out a decrease in the next event.

\section{From probability to repeated computer experiments}

Suppose we repeat a complete simulation many times, always starting with
the same parameters and initial count. One complete run is a
\textbf{replicate}. Each replicate ends with A either lost or fixed.

Let \(R\) denote the number of replicates. Let \(K\) denote the number of
those replicates in which A fixes. A natural estimate of the fixation
probability is
\[
 \widehat p=\frac{K}{R}.
\]
The symbol \(\widehat p\), read as ``p hat'', means an estimated probability.
The hat reminds us that this number is computed from a finite sample of
random outcomes. We use the letter \(p\), without a hat, for the underlying
probability specified by the model.

If twelve out of one hundred runs fix, then \(\widehat p=12/100=0.12\).
Another batch might give nine or fourteen fixations. The rules can stay the
same while the observed count differs.

Larger numbers of replicates usually give more precise estimates. They do
not change the underlying model probability. A later chapter will put a
confidence interval around the estimate and explain what that interval means.

A \textbf{random seed} is an initial setting for the computer's random-number
generator. Recording a seed helps reproduce a particular computation. It
does not turn one random outcome into a universal result.

\begin{checkpoint}
\textbf{Question 1.} A parent draw selects A with probability \(0.4\).
An independent replacement draw selects B with probability \(0.6\).
What is the probability of both outcomes?

\textbf{Answer.} Multiply the probabilities:
\(0.4\times0.6=0.24\). We multiply because both outcomes are required,
and the two draws are independent given this state.

\textbf{Question 2.} The count increases with probability \(0.24\) and
decreases with probability \(0.24\). What is the probability of no change?
What is the expected count change?

\textbf{Answer.} The probability of a change is \(0.24+0.24=0.48\).
The probability of no change is \(1-0.48=0.52\). The expected change is
\(1(0.24)+(-1)(0.24)+0(0.52)=0\). An expected change of zero does not
mean that every event leaves the count unchanged.

\textbf{Question 3.} A set of \(R=200\) complete runs contains \(K=30\)
fixations. What probability do we estimate?

\textbf{Answer.} The estimate is
\(\widehat p=K/R=30/200=0.15\). Fifteen percent of this particular batch
fixed. The true model probability need not be exactly \(0.15\).
\end{checkpoint}

\section{The three questions to ask before a calculation}

Before multiplying, ask whether you need \emph{both} events and whether the
second probability is conditional on the first. Before adding, ask whether
the alternatives overlap. Before interpreting an average, ask whether it is
a statement about many repetitions or about one actual trajectory.

Those checks will carry most of the probability reasoning in this project.
